\documentclass[11pt]{article}
% ======================================================================
%  examstyle.tex — EPFL-style exam layout for the weekly Time Series exams
%  Shared by every Exam_Week_NN(.tex / _Solutions.tex). Compiles with pdflatex.
% ======================================================================
\usepackage[utf8]{inputenc}
\usepackage[T1]{fontenc}
\usepackage[english]{babel}
\usepackage[a4paper,margin=2.4cm]{geometry}
\usepackage{amsmath,amssymb,amsthm,mathtools}
\usepackage{bm}
\usepackage{enumitem}
\usepackage{array}

\setlength{\parindent}{0pt}
\setlength{\parskip}{0.5em}
\emergencystretch=3em
\allowdisplaybreaks[1]

% ---------- math macros (same names as the rest of the project) ----------
\newcommand{\E}{\mathbb{E}}
\DeclareMathOperator{\Var}{Var}
\DeclareMathOperator{\Cov}{Cov}
\DeclareMathOperator{\Corr}{Corr}
\DeclareMathOperator*{\argmin}{arg\,min}
\DeclareMathOperator{\MSE}{MSE}
\DeclareMathOperator{\trace}{tr}
\newcommand{\Reals}{\mathbb{R}}
\newcommand{\Ints}{\mathbb{Z}}
\newcommand{\Comp}{\mathbb{C}}
\newcommand{\Nat}{\mathbb{N}}
\newcommand{\Prob}{\mathbb{P}}
\newcommand{\Tr}{^{\mathsf{T}}}
\newcommand{\conj}[1]{\overline{#1}}
\newcommand{\abs}[1]{\left\lvert #1 \right\rvert}
\newcommand{\norm}[1]{\left\lVert #1 \right\rVert}
\newcommand{\ind}{\mathbf{1}}
\newcommand{\eps}{\varepsilon}
\newcommand{\diff}{\nabla}
\newcommand{\acvs}{\gamma}
\newcommand{\acf}{\rho}
\newcommand{\pacf}{\alpha}
\newcommand{\sdf}{\mathcal{S}}
\newcommand{\Bsh}{B}
\newcommand{\iid}{\overset{\text{iid}}{\sim}}

\setlist[enumerate]{leftmargin=2em,topsep=2pt,itemsep=4pt}
\setlist[itemize]{leftmargin=1.6em,topsep=2pt,itemsep=2pt}

\newcounter{exo}

% ---------- exam cover header :  \examheader{exam title}{duration} ----------
\newcommand{\examheader}[2]{%
  \thispagestyle{empty}
  \begin{center}
    {\Large\bfseries Time Series}\\[3pt]
    Spring Semester 2026, EPFL\\
    \'Ecole Polytechnique F\'ed\'erale de Lausanne\\
    Prof.\ Dr.\ Sofia C.\ Olhede\\[7pt]
    \hrule height 0.8pt
    \vspace{9pt}
    {\large\bfseries Practice Exam --- #1}\\[4pt]
    {\normalsize Duration: #2 \quad$\cdot$\quad No calculator \quad$\cdot$\quad Answer on paper}\\
    \vspace{8pt}
    \hrule height 0.8pt
  \end{center}
  \vspace{14pt}
  \begin{tabular}{@{}p{8.2cm}@{}p{8cm}@{}}
    Name:\enspace\hrulefill & Surname:\enspace\hrulefill
  \end{tabular}\\[14pt]
  SCIPER (Student Card Number):\enspace\hrulefill
  \vspace{16pt}
}

% ---------- points table (fixed: 4 problems) ----------
\newcommand{\pointstable}{%
  \begin{center}
  \renewcommand{\arraystretch}{1.5}
  \begin{tabular}{|p{5cm}|p{4cm}|}
    \hline
    \textbf{Exercise} & \textbf{Points}\\\hline
    1 & \\\hline
    2 & \\\hline
    3 & \\\hline
    4 & \\\hline
    \textbf{Total} & \\\hline
  \end{tabular}
  \end{center}
}

% ---------- problem heading (exam: one problem per page) ----------
\newcommand{\problem}[1]{%
  \clearpage
  \stepcounter{exo}%
  \begin{center}\textbf{\large Problem \theexo\ (#1 points)}\end{center}
  \vspace{4pt}
}
% ---------- answer space: fill the statement page + one blank answer page ----------
% Put \answerpage at the END of each problem so every exercise gets its own page(s).
\newcommand{\answerpage}{\par\vfill\newpage\null\newpage}

% ---------- solutions document ----------
\newcommand{\solheader}[1]{%
  \begin{center}
    {\Large\bfseries Time Series --- Solutions}\\[3pt]
    {\large Practice Exam --- #1}\\[2pt]
    EPFL, MATH-342 \quad$\cdot$\quad Prof.\ S.\ C.\ Olhede\\[5pt]
    \hrule height 0.8pt
  \end{center}
  \vspace{10pt}
}
\newcommand{\solproblem}[1]{%
  \bigskip
  \stepcounter{exo}%
  {\large\bfseries Problem \theexo\ (#1 points)}\par\nobreak\vspace{2pt}
}
% Rappel de l'énoncé avant la solution (dans les corrigés)
\newcommand{\statementstart}{\par\vspace{1pt}{\bfseries\itshape Statement.}\par\nobreak\begingroup\small\itshape}
\newcommand{\statementend}{\par\endgroup\vspace{5pt}\hrule\vspace{6pt}{\bfseries Solution.}\par\nobreak\vspace{2pt}}

\begin{document}

\solheader{Week 5: Trends, Seasonality \& SARIMA}

% ---------------------------------------------------------------
\solproblem{5}
\statementstart
Consider the signal-plus-noise model $X_t=\mu_t+Y_t$, where $\mu_t$ is a deterministic trend and $\{Y_t\}$ is a zero-mean stationary process with autocovariance sequence $\acvs_\tau$. Recall the difference operator $\diff=I-\Bsh$ and the seasonal difference $\diff_{(s)}=I-\Bsh^{s}$.

\textbf{(a)} Take $\mu_t=a+bt$ (linear trend). Show that $\diff X_t=b+\diff Y_t$, and explain in one sentence why $\diff Y_t$ is stationary while $\{X_t\}$ in general is not. \hfill(1.5 pts)

\textbf{(b)} Now let $\mu_t=a+bt+s_t$ with a \emph{quarterly} season $s_{t+4}=s_t$. Compute $\diff\diff_{(4)}X_t$ \emph{explicitly} as a linear combination of the $Y$'s, and verify that its mean is $0$ and that it is weakly stationary. \hfill(2 pts)

\textbf{(c)} \emph{(A classic trap.)} The operators $\diff_{(s)}=I-\Bsh^{s}$ and $\diff^{s}=(I-\Bsh)^{s}$ are \emph{not} the same. For the deterministic sequence $\mu_t=t^{2}$, compute both $\diff^{2}\mu_t$ and $\diff_{(2)}\mu_t$ and state which kind of structure each operator removes. \hfill(1.5 pts)
\statementend

\textbf{(a)} With $\mu_t=a+bt$, linearity of $\diff=I-\Bsh$ gives
\[
\diff X_t=X_t-X_{t-1}
=\bigl(a+bt+Y_t\bigr)-\bigl(a+b(t-1)+Y_{t-1}\bigr)
=b+\bigl(Y_t-Y_{t-1}\bigr)=b+\diff Y_t .
\]
So differencing collapses the linear trend $a+bt$ to the constant $b$. Now $\diff Y_t=Y_t-Y_{t-1}$ is a \emph{fixed finite linear combination} of a stationary process, hence itself stationary (its mean is $0$ and its covariance depends on the lag only). The original $\{X_t\}$ is not stationary because $\E[X_t]=a+bt$ depends on $t$.

\textbf{(b)} With $\mu_t=a+bt+s_t$ and $s_{t+4}=s_t$, apply the seasonal difference first. Using $s_t=s_{t-4}$,
\[
\diff_{(4)}X_t=X_t-X_{t-4}
=\bigl(a+bt+s_t+Y_t\bigr)-\bigl(a+b(t-4)+s_{t-4}+Y_{t-4}\bigr)
=4b+\bigl(Y_t-Y_{t-4}\bigr).
\]
The period-$4$ season has been annihilated and the linear trend has collapsed to the constant $4b$. Now apply $\diff=I-\Bsh$, which kills the constant:
\[
\diff\diff_{(4)}X_t
=\bigl(4b+Y_t-Y_{t-4}\bigr)-\bigl(4b+Y_{t-1}-Y_{t-5}\bigr)
=\boxed{\,Y_t-Y_{t-1}-Y_{t-4}+Y_{t-5}\,}.
\]
This is a fixed linear combination of $\{Y_t\}$, so its mean is $\E[Y_t-Y_{t-1}-Y_{t-4}+Y_{t-5}]=0$, and by bilinearity its covariance,
\[
\begin{aligned}
\Cov\bigl(\diff\diff_{(4)}X_t,\ \diff\diff_{(4)}X_{t-\tau}\bigr)
=\;&4\acvs_\tau
-2\acvs_{\tau-1}-2\acvs_{\tau+1}
-2\acvs_{\tau-4}-2\acvs_{\tau+4}\\
&+\acvs_{\tau-3}+\acvs_{\tau+3}
+\acvs_{\tau-5}+\acvs_{\tau+5}
+\acvs_{\tau+1-4}+\acvs_{\tau-1+4},
\end{aligned}
\]
depends on the lag $\tau$ only (no $t$). Hence $\diff\diff_{(4)}X_t$ is weakly stationary; all $t$-dependence has been removed. (One need not simplify the covariance fully; the key point is that it is a function of $\tau$ alone.)

\textbf{(c)} Treat $\mu_t=t^{2}$ as a deterministic sequence.

\emph{Ordinary second difference $\diff^{2}=(I-\Bsh)^{2}$:}
\[
\diff^{2}t^{2}=t^{2}-2(t-1)^{2}+(t-2)^{2}
=t^{2}-2\bigl(t^{2}-2t+1\bigr)+\bigl(t^{2}-4t+4\bigr)
=2 ,
\]
a constant. (And $\diff^{3}t^{2}=\diff\,(2)=0$.) So $\diff^{q}$ \emph{annihilates a polynomial trend of degree $q-1$}: here $\diff^{2}$ reduces the degree-$2$ polynomial $t^{2}$ to a degree-$0$ constant, and one further $\diff$ kills it.

\emph{Seasonal difference of period $2$, $\diff_{(2)}=I-\Bsh^{2}$:}
\[
\diff_{(2)}t^{2}=t^{2}-(t-2)^{2}=t^{2}-\bigl(t^{2}-4t+4\bigr)=4t-4 ,
\]
which is still linear (not annihilated). The seasonal operator $\diff_{(s)}$ removes a deterministic \emph{period-$s$ season} ($\diff_{(s)}s_t=s_t-s_{t-s}=0$), \emph{not} a polynomial trend. The two operators serve different purposes, and $\diff_{(2)}\neq\diff^{2}$.

% ---------------------------------------------------------------
\solproblem{5}
\statementstart
Let $\{X_t\}$ be the multiplicative seasonal process $\mathrm{SARMA}(0,1)\times(1,0)_{4}$ (quarterly data), defined by
\[
\bigl(1-\Phi\,\Bsh^{4}\bigr)X_t=\bigl(1-\theta\Bsh\bigr)\eps_t,
\qquad \abs\Phi<1 ,
\]
a seasonal autoregression at lag $4$ driven by a non-seasonal $\mathrm{MA}(1)$ input.

\textbf{(a)} Write the explicit one-step recursion for $X_t$, and by inverting the seasonal AR factor obtain the $\mathrm{MA}(\infty)$ weights $\psi_j$ in $X_t=\sum_{j\ge0}\psi_j\eps_{t-j}$ (give a closed form for $\psi_{4l}$ and $\psi_{4l+1}$, and state that all other $\psi_j=0$). \hfill(2 pts)

\textbf{(b)} Using the white-noise filter rule $\acvs_\tau=\sigma^{2}\sum_{j\ge0}\psi_j\psi_{j+\tau}$, compute $\acvs_0,\acvs_1,\acvs_3,\acvs_4,\acvs_5$ in closed form, and explain why the autocovariance lives in clusters around the seasonal lags $0,4,8,\dots$ \hfill(2 pts)

\textbf{(c)} Evaluate the autocorrelations $\acf_1,\acf_3,\acf_4$ numerically for $\Phi=\tfrac12$, $\theta=\tfrac25$, $\sigma^{2}=1$. \hfill(1 pt)
\statementend

\textbf{(a)} Expand the seasonal AR factor $(1-\Phi\Bsh^{4})X_t=(1-\theta\Bsh)\eps_t$ and solve for $X_t$:
\[
\boxed{\,X_t=\Phi\,X_{t-4}+\eps_t-\theta\eps_{t-1}\,}
\]
--- an $\mathrm{AR}(1)$ acting only on the seasonal lag $4$, driven by an $\mathrm{MA}(1)$ input. Invert the seasonal factor as a geometric series in $\Bsh^{4}$ (valid since $\abs\Phi<1$),
\[
\frac{1}{1-\Phi\Bsh^{4}}=\sum_{l\ge0}\Phi^{l}\Bsh^{4l},
\]
so that
\[
X_t=\sum_{l\ge0}\Phi^{l}\Bsh^{4l}\bigl(\eps_t-\theta\eps_{t-1}\bigr)
=\sum_{l\ge0}\Phi^{l}\bigl(\eps_{t-4l}-\theta\eps_{t-4l-1}\bigr).
\]
Reading off the coefficient of $\eps_{t-j}$:
\[
\boxed{\;\psi_{4l}=\Phi^{l},\qquad \psi_{4l+1}=-\theta\,\Phi^{l}\quad(l\ge0),\qquad \psi_j=0\ \text{otherwise.}\;}
\]
(Check: $\psi_0=1,\ \psi_1=-\theta,\ \psi_4=\Phi,\ \psi_5=-\theta\Phi,\dots$, and indeed $\psi_j-\Phi\psi_{j-4}$ reproduces the $\mathrm{MA}(1)$ numerator $1-\theta\Bsh$.)

\textbf{(b)} Apply $\acvs_\tau=\sigma^{2}\sum_{j\ge0}\psi_j\psi_{j+\tau}$. The only nonzero taps sit at $j\equiv0$ and $j\equiv1\pmod4$, so a product $\psi_j\psi_{j+\tau}$ survives only when both indices are of this type; this happens at $\tau\in\{0,\pm1,\pm3,\pm4,\pm5,\dots\}$, i.e.\ in clusters $\{4k-1,\,4k,\,4k+1\}$ around the seasonal multiples $4k$.

\emph{Lag $0$:} $\sum_l\psi_{4l}^{2}+\sum_l\psi_{4l+1}^{2}=\sum_l\Phi^{2l}+\theta^{2}\sum_l\Phi^{2l}=(1+\theta^{2})/(1-\Phi^{2})$, so
\[
\acvs_0=\frac{(1+\theta^{2})\,\sigma^{2}}{1-\Phi^{2}} .
\]
\emph{Lag $1$:} the surviving pairs are $\psi_{4l}\psi_{4l+1}=\Phi^{l}(-\theta\Phi^{l})=-\theta\Phi^{2l}$, so
\[
\acvs_1=\sigma^{2}\sum_{l\ge0}\bigl(-\theta\Phi^{2l}\bigr)=\frac{-\theta\,\sigma^{2}}{1-\Phi^{2}} .
\]
\emph{Lag $3$:} pairs $\psi_{4l+1}\psi_{4l+4}=(-\theta\Phi^{l})(\Phi^{l+1})=-\theta\Phi^{2l+1}$, so
\[
\acvs_3=\sigma^{2}\sum_{l\ge0}\bigl(-\theta\Phi^{2l+1}\bigr)=\frac{-\theta\,\Phi\,\sigma^{2}}{1-\Phi^{2}} .
\]
\emph{Lag $4$:} pairs $\psi_{4l}\psi_{4l+4}=\Phi^{2l+1}$ and $\psi_{4l+1}\psi_{4l+5}=(-\theta\Phi^{l})(-\theta\Phi^{l+1})=\theta^{2}\Phi^{2l+1}$, so
\[
\acvs_4=\sigma^{2}\sum_{l\ge0}(1+\theta^{2})\Phi^{2l+1}=\frac{(1+\theta^{2})\Phi\,\sigma^{2}}{1-\Phi^{2}}=\Phi\,\acvs_0 .
\]
\emph{Lag $5$:} pairs $\psi_{4l}\psi_{4l+5}=\Phi^{l}(-\theta\Phi^{l+1})=-\theta\Phi^{2l+1}$, so
\[
\acvs_5=\frac{-\theta\,\Phi\,\sigma^{2}}{1-\Phi^{2}}=\acvs_3 .
\]
Collecting,
\[
\boxed{\;
\acvs_0=\frac{(1+\theta^{2})\sigma^{2}}{1-\Phi^{2}},\quad
\acvs_1=\frac{-\theta\,\sigma^{2}}{1-\Phi^{2}},\quad
\acvs_3=\acvs_5=\frac{-\theta\Phi\,\sigma^{2}}{1-\Phi^{2}},\quad
\acvs_4=\Phi\,\acvs_0 .
\;}
\]
So the seasonal $\mathrm{AR}(1)$ produces a spike at every seasonal lag $4k$ of size $\Phi^{k}\acvs_0$, each flanked by $\mathrm{MA}(1)$ side-lobes at $4k\pm1$; the whole pattern decays geometrically in $\Phi$.

\textbf{(c)} With $\Phi=\tfrac12$, $\theta=\tfrac25$, $\sigma^{2}=1$: $1-\Phi^{2}=\tfrac34$ and $1+\theta^{2}=1+\tfrac{4}{25}=\tfrac{29}{25}$. Hence
\[
\acvs_0=\frac{29/25}{3/4}=\frac{116}{75}\approx1.5467 .
\]
Then
\[
\acf_1=\frac{\acvs_1}{\acvs_0}=\frac{-\theta}{1+\theta^{2}}=\frac{-2/5}{29/25}=\frac{-2/5\cdot25}{29}=-\frac{10}{29}\approx-0.3448 ,
\]
\[
\acf_4=\frac{\acvs_4}{\acvs_0}=\Phi=\tfrac12=0.5 ,
\qquad
\acf_3=\frac{\acvs_3}{\acvs_0}=\Phi\,\acf_1=\tfrac12\cdot\Bigl(-\tfrac{10}{29}\Bigr)=-\frac{5}{29}\approx-0.1724 .
\]
(So the correlogram shows a dominant seasonal spike $\acf_4=0.5$ with small negative side-lobes $\acf_3=\acf_5\approx-0.172$ at $4\pm1$, and $\acf_1\approx-0.345$ near the origin.)

% ---------------------------------------------------------------
\solproblem{5}
\statementstart
This problem studies non-stationary integrated models.

\textbf{(a)} \emph{(Random walk.)} Let $X_t=\sum_{j=1}^{t}\eps_j$ with $X_0=0$ (an $\mathrm{ARIMA}(0,1,0)$). Compute $\E[X_t]$, $\Var(X_t)$ and $\Cov(X_t,X_{t+\tau})$ for $\tau\ge0$, deduce that $\{X_t\}$ is \emph{not} stationary, and show $\diff X_t$ is white noise. Show also that $\Corr(X_t,X_{t+\tau})\to1$ as $t\to\infty$ with $\tau$ fixed. \hfill(2 pts)

\textbf{(b)} \emph{(IMA$(1,1)$.)} Let $Y_t=Y_{t-1}+\eps_t-\theta\eps_{t-1}$ with $Y_0=0$. By unrolling the recursion, show that for $t\ge1$
\[
Y_t=\eps_t+(1-\theta)\sum_{j=1}^{t-1}\eps_j-\theta\eps_0 ,
\]
and hence compute $\Var(Y_t)$. For $\theta\neq1$, what is the order of growth of $\Var(Y_t)$ in $t$? \hfill(2 pts)

\textbf{(c)} \emph{(Over-differencing.)} What happens to $Y_t$ and to $\Var(Y_t)$ at the boundary $\theta=1$? Show that $\diff Y_t$ is then a \emph{non-invertible} $\mathrm{MA}(1)$, and state the value of its lag-$1$ autocorrelation $\acf_1$. Why is this the warning signature of over-differencing? \hfill(1 pt)
\statementend

\textbf{(a)} Since $X_t=\sum_{j=1}^{t}\eps_j$ is a sum of mean-zero noise, $\E[X_t]=0$. As the $\eps_j$ are uncorrelated with variance $\sigma^{2}$,
\[
\Var(X_t)=\sum_{j=1}^{t}\Var(\eps_j)=t\,\sigma^{2}.
\]
For $\tau\ge0$, $X_{t+\tau}=X_t+\sum_{j=t+1}^{t+\tau}\eps_j$ adds only \emph{new} (uncorrelated) noise, so the overlap is the first $t$ shared terms:
\[
\Cov(X_t,X_{t+\tau})=\Cov\Bigl(\sum_{j=1}^{t}\eps_j,\ \sum_{k=1}^{t+\tau}\eps_k\Bigr)
=\sum_{j=1}^{t}\sigma^{2}=t\,\sigma^{2}.
\]
Both $\Var(X_t)$ and $\Cov(X_t,X_{t+\tau})$ depend on $t$, so $\{X_t\}$ is \textbf{not} (weakly) stationary. However
\[
\diff X_t=X_t-X_{t-1}=\eps_t ,
\]
which is white noise; hence $\{X_t\}$ is an $\mathrm{ARIMA}(0,1,0)$. Finally
\[
\Corr(X_t,X_{t+\tau})=\frac{t\sigma^{2}}{\sqrt{t\sigma^{2}}\sqrt{(t+\tau)\sigma^{2}}}
=\sqrt{\frac{t}{t+\tau}}\xrightarrow[t\to\infty]{}1 \quad(\tau\text{ fixed}),
\]
the near-unit correlation of nearby values characteristic of an integrated series.

\textbf{(b)} The recursion $Y_t=Y_{t-1}+(\eps_t-\theta\eps_{t-1})$ with $Y_0=0$ telescopes:
\[
Y_t=\sum_{k=1}^{t}\bigl(\eps_k-\theta\eps_{k-1}\bigr)
=\sum_{k=1}^{t}\eps_k-\theta\sum_{k=1}^{t}\eps_{k-1}
=\sum_{k=1}^{t}\eps_k-\theta\sum_{j=0}^{t-1}\eps_j .
\]
Separating the unmatched endpoints $\eps_t$ (only in the first sum) and $\eps_0$ (only in the second),
\[
Y_t=\eps_t+\sum_{k=1}^{t-1}\eps_k-\theta\eps_0-\theta\sum_{j=1}^{t-1}\eps_j
=\eps_t+(1-\theta)\sum_{j=1}^{t-1}\eps_j-\theta\eps_0 ,
\]
as claimed. The coefficients of the $t+1$ independent innovations $\eps_0,\eps_1,\dots,\eps_t$ are: $-\theta$ on $\eps_0$, $(1-\theta)$ on each of $\eps_1,\dots,\eps_{t-1}$ (that is $t-1$ terms), and $1$ on $\eps_t$. By independence,
\[
\boxed{\;\Var(Y_t)=\sigma^{2}\Bigl[\,1+(1-\theta)^{2}(t-1)+\theta^{2}\,\Bigr].\;}
\]
For $\theta\neq1$ this is \emph{linear in $t$} with positive slope $(1-\theta)^{2}\sigma^{2}>0$: $\Var(Y_t)=O(t)$, exactly the random-walk-like variance growth that flags an integrated (needs-differencing) series.

\textbf{(c)} At the boundary $\theta=1$ the slope $(1-\theta)^{2}$ vanishes, so
\[
\Var(Y_t)=\sigma^{2}\bigl[1+0+1\bigr]=2\sigma^{2},
\]
which is \emph{bounded} (no growth). Indeed the recursion telescopes completely:
\[
Y_t=\sum_{k=1}^{t}(\eps_k-\eps_{k-1})=\eps_t-\eps_0 ,
\]
so the level is itself stationary, and
\[
\diff Y_t=Y_t-Y_{t-1}=\eps_t-\eps_{t-1},
\]
an $\mathrm{MA}(1)$ with unit coefficient $\theta=1$. Its MA polynomial $1-z$ has root $z=1$ \emph{on} the unit circle, so it is \textbf{non-invertible}, and its lag-$1$ autocorrelation is
\[
\acf_1=\frac{-\theta}{1+\theta^{2}}=\frac{-1}{1+1}=-\tfrac12 .
\]
This is the warning signature of \textbf{over-differencing}: differencing a series that was already stationary injects a $(1-\Bsh)$ factor into the MA part, i.e.\ a non-invertible $\mathrm{MA}(1)$ with $\acf_1=-\tfrac12$. Seeing a sample $\acf_1\approx-\tfrac12$ after a difference is the cue to difference one time fewer.

% ---------------------------------------------------------------
\solproblem{5}
\statementstart
\textbf{(Revision of Week 4: Yule--Walker estimation.)} Let $\{X_t\}$ be a mean-zero causal $\mathrm{AR}(2)$ process $X_t=\phi_1X_{t-1}+\phi_2X_{t-2}+\eps_t$ observed at $X_1,\dots,X_n$, with sample autocovariances
\[
\hat\acvs_0=10,\qquad \hat\acvs_1=5,\qquad \hat\acvs_2=1 .
\]

\textbf{(a)} Derive the Yule--Walker equations $\acvs_k=\phi_1\acvs_{k-1}+\phi_2\acvs_{k-2}$ ($k\ge1$) and the variance identity $\acvs_0=\phi_1\acvs_1+\phi_2\acvs_2+\sigma^{2}$, making clear where \emph{causality} is used. \hfill(2 pts)

\textbf{(b)} Set up the $2\times2$ Yule--Walker system, invert it \emph{by hand}, and report the estimates $\hat\phi_1,\hat\phi_2$ and $\hat\sigma^{2}$. \hfill(2 pts)

\textbf{(c)} Verify that the fitted model is causal (stationary), e.g.\ by locating the roots of $\hat\Phi(z)$ or via the $\mathrm{AR}(2)$ stationarity triangle. \hfill(1 pt)
\statementend

\textbf{(a)} Multiply the model $X_t=\phi_1X_{t-1}+\phi_2X_{t-2}+\eps_t$ by $X_{t-k}$ and take expectations. The process is causal, so $X_{t-k}$ is a one-sided combination of $\eps_{t-k},\eps_{t-k-1},\dots$; hence for $k\ge1$ the noise term is uncorrelated with the past, $\E[\eps_tX_{t-k}]=0$. With mean zero, $\E[X_{t-j}X_{t-k}]=\acvs_{k-j}$, giving for $k\ge1$
\[
\acvs_k=\phi_1\acvs_{k-1}+\phi_2\acvs_{k-2}\qquad(\text{using }\acvs_{-1}=\acvs_1).
\]
Multiplying instead by $X_t$ keeps the noise term: by causality and $\psi_0=1$,
\[
\E[\eps_tX_t]=\E\bigl[\eps_t(\phi_1X_{t-1}+\phi_2X_{t-2}+\eps_t)\bigr]=\sigma^{2},
\]
so the variance identity is
\[
\acvs_0=\phi_1\acvs_1+\phi_2\acvs_2+\sigma^{2}.
\]
The single fact ``causality $\Rightarrow\E[\eps_tX_{t-k}]=0$ for $k\ge1$'' is what closes the system. Replacing the population $\acvs$ by the sample $\hat\acvs$ gives the Yule--Walker \emph{estimators}.

\textbf{(b)} Taking $k=1,2$ gives the $2\times2$ Yule--Walker system
$\hat\Gamma_2\hat{\boldsymbol\phi}=\hat{\boldsymbol\acvs}_2$:
\[
\begin{pmatrix}\hat\acvs_0 & \hat\acvs_1\\[2pt] \hat\acvs_1 & \hat\acvs_0\end{pmatrix}
\begin{pmatrix}\hat\phi_1\\[2pt] \hat\phi_2\end{pmatrix}
=\begin{pmatrix}\hat\acvs_1\\[2pt] \hat\acvs_2\end{pmatrix}
\quad\Longrightarrow\quad
\begin{pmatrix}10 & 5\\[2pt] 5 & 10\end{pmatrix}
\begin{pmatrix}\hat\phi_1\\[2pt] \hat\phi_2\end{pmatrix}
=\begin{pmatrix}5\\[2pt] 1\end{pmatrix}.
\]
The inverse of $\left(\begin{smallmatrix}a&b\\ b&a\end{smallmatrix}\right)$ is $\tfrac{1}{a^{2}-b^{2}}\left(\begin{smallmatrix}a&-b\\ -b&a\end{smallmatrix}\right)$; here $a=10,\ b=5$ and $a^{2}-b^{2}=100-25=75$. Hence
\[
\begin{pmatrix}\hat\phi_1\\ \hat\phi_2\end{pmatrix}
=\frac{1}{75}\begin{pmatrix}10 & -5\\ -5 & 10\end{pmatrix}\begin{pmatrix}5\\ 1\end{pmatrix}
=\frac{1}{75}\begin{pmatrix}10\cdot5-5\cdot1\\ -5\cdot5+10\cdot1\end{pmatrix}
=\frac{1}{75}\begin{pmatrix}45\\ -15\end{pmatrix}
=\begin{pmatrix}0.6\\ -0.2\end{pmatrix}.
\]
The innovation variance is
\[
\hat\sigma^{2}=\hat\acvs_0-\hat\phi_1\hat\acvs_1-\hat\phi_2\hat\acvs_2
=10-0.6\cdot5-(-0.2)\cdot1=10-3+0.2=7.2 .
\]
So $\boxed{\hat\phi_1=0.6,\ \hat\phi_2=-0.2,\ \hat\sigma^{2}=7.2.}$
(Consistency check via $\acvs_0=\sigma^{2}/(1-\phi_1\acf_1-\phi_2\acf_2)$ with $\acf_1=\tfrac12,\acf_2=\tfrac{1}{10}$: $1-0.6\cdot\tfrac12-(-0.2)\cdot\tfrac1{10}=1-0.3+0.02=0.72$, and $7.2/0.72=10=\hat\acvs_0$.\ \checkmark)

\textbf{(c)} The fitted AR polynomial is $\hat\Phi(z)=1-0.6z+0.2z^{2}$. Its roots solve $0.2z^{2}-0.6z+1=0$, i.e.\ $z^{2}-3z+5=0$:
\[
z=\frac{3\pm\sqrt{9-20}}{2}=\frac{3\pm\sqrt{-11}}{2}=\tfrac32\pm\tfrac{\sqrt{11}}{2}\,i,
\qquad
\abs{z}=\sqrt{\Bigl(\tfrac32\Bigr)^{2}+\tfrac{11}{4}}=\sqrt{\tfrac{9}{4}+\tfrac{11}{4}}=\sqrt5\approx2.236 .
\]
Both roots have modulus $\sqrt5>1$ (outside the unit circle), so the fit is \textbf{causal/stationary}. Equivalently, the $\mathrm{AR}(2)$ stationarity-triangle conditions hold: $\hat\phi_1+\hat\phi_2=0.4<1$, $\hat\phi_2-\hat\phi_1=-0.8<1$, and $\abs{\hat\phi_2}=0.2<1$. (Unlike least squares, the Yule--Walker estimator always returns a stationary fit.)

\end{document}
